%%%%%%%%%%  scan idx 203  |  folio 190  %%%%%%%%%%
\SGAPrintPage{190}%
as are the inclusions $f_{i_0\ldots i_p}:U_{i_0\ldots i_p}=U_{i_0}\cap\ldots\cap U_{i_p}\longrightarrow S$. Let $E\in\Ob\Qcoh(S)$, and consider the resolution of $E$ defined by the covering $(U_i)$ (T.F. II 5.2):
\sgaStableTarget{sga6.e02.eq.unnamed.occ3}\begin{equation}\tag{*}
 0\longrightarrow E\longrightarrow \coprod f_{i*}f_i^*(E)\longrightarrow\ldots\longrightarrow
 \coprod f_{i_0\ldots i_p*}f_{i_0\ldots i_p}^*(E)\longrightarrow .
\end{equation}
It gives a spectral sequence
\sgaStableTarget{sga6.e02.eq.unnamed.occ4}\begin{equation}\tag{**}
 E_1^{pq}=\coprod_{i_0\ldots i_p}R^qQ\bigl(f_{i_0\ldots i_p*}f_{i_0\ldots i_p}^*(E)\bigr)
 \Longrightarrow R^{*}Q(E).
\end{equation}
Since $f_{i_0\ldots i_p}$ is affine, (3.4) gives:
\sgaStableTarget{sga6.e02.eq.unnamed.occ5}\begin{equation}\tag{***}
 R^qQ\bigl(f_{i_0\ldots i_p*}f_{i_0\ldots i_p}^*(E)\bigr)
 \simeq f_{i_0\ldots i_p*}R^qQ\bigl(f_{i_0\ldots i_p}^*(E)\bigr)
\end{equation}
for every $q$, whence in particular
\[
 E_1^{pq}=0\quad\text{for }q>0,
\]
by the affine case, already treated. From (**) we deduce isomorphisms
\[
 R^nQ(E)\simeq H^n(QE^\bullet)\quad\text{for }n\geq0,
\]
where $E^\bullet$ is the resolution of $E$ defined by the exact sequence (*). But (***) gives $QE^n\simeq E^n$ for $n\geq0$, whence finally
\[
 R^nQ(E)=0\quad\text{for }n>0,
\]
and $Q(E)\simeq E$, q.e.d.

\uline{Proof of} (3.5.2). The spectral sequence
\[
 R^pQ(H^q(E))\Longrightarrow R^{*}Q(E),
\]
%%%%%%%%%%  scan idx 204  |  folio 191  %%%%%%%%%%
\SGAPrintPage{191}%
immediately reduces us to the case where $E\in\Ob\Qcoh(S)$. Since the composite
\[
 E\longrightarrow \underline{\underline{R}}QE\longrightarrow E
\]
of the adjunction arrows is the identity, the assertion therefore follows from (3.5.1). This completes the proof of (3.5).

\sgaStableTarget{sga6.e02.rem.3-6}\textbf{Remark 3.6.} The conclusion of (3.5) would be false with the hypothesis \textquotedblleft separated or noetherian\textquotedblright{} omitted, as a neat counterexample of J.L. Verdier shows (see Appendix I to this exposé).

\sgaStableTarget{sga6.e02.prop.3-7}\textbf{Proposition 3.7.} \uline{Under the hypothesis of} (3.5)\uline{, suppose further that there exists an integer} $N$ \uline{such that, for every} $F\in\Ob\Mod(S)$ \uline{and every quasi-compact open subset} $U$ \uline{of} $S$,
\[
 H^i(U,F)=0\quad\text{for }i>N\qquad\text{\SGAAuthorFootMark{1}}\ .
\]
\SGAAuthorFootText{1}{By (Tohoku 3.6.5), this condition holds when the underlying space of $S$ is noetherian and has Krull dimension $\leq N$.}
\uline{Then:}
\begin{enumerate}[label=\alph*)]
\item \uline{The coherer has finite cohomological dimension.}
\item \uline{The functor}
\[
 \phi:\mathrm D(\Qcoh(S))\longrightarrow \mathrm D(S)
\]
\uline{is fully faithful and has as essential image the full subcategory} $\mathrm D(S)_{\qcoh}$ \uline{of} $\mathrm D(S)$ \uline{formed by complexes with quasi-coherent cohomology.}
\end{enumerate}

\uline{Proof of} a). When $S$ is affine, the formula $Q=\widetilde{\Gamma}_S$ (cf. proof of (3.2)) implies that $Q$ has cohomological dimension $\leq N$. We pass to the general case. Let us show that there exists an integer $a$ such that $R^qQ(E)=0$ for every $E\in\Ob\Mod(S)$ and every $q>a$. The hypothesis implies that every $E\in\Ob\Mod(S)$ admits a finite right resolution $E^\bullet$ of length $N$ such that the $E^i$ are acyclic, i.e.\ satisfy $H^q(U,E^i)=0$ for every open subset $U$ of $S$ and every $q>0$. It therefore suffices to prove that there exists an integer $b$ such that $R^qQ(E)=0$ for every acyclic $\underline{O}_S$-Module $E$ and every $q>b$. (We may then take $a=b+N$). Let $(U_i)$ be a finite covering of $S$ by $N'$ affine
%%%%%%%%%%  scan idx 205  |  folio 192  %%%%%%%%%%
\SGAPrintPage{192}%
open subsets. Let $E$ be an acyclic $\underline{O}_S$-Module, and consider the resolution $C^\bullet(E)$ of $E$ defined by the covering $(U_i)$ (T.F. II 5.2). There are two cases to examine:

1°) $S$ is \uline{separated}. The components of $C^\bullet(E)$ are then sums of terms of the form $f_*f^*(E)$, where $f:U\longrightarrow S$ is the inclusion of an affine open subset (hence an affine morphism, since $S$ is separated). We have
\[
 \underline{\underline{R}}(Qf_*)f^*E\simeq \underline{\underline{R}}Q\,\underline{\underline{R}}f_*f^*E\simeq \underline{\underline{R}}Q\,f_*f^*E,
\]
since $f^*E$ is acyclic. On the other hand, by (3.4), we have
\[
 \underline{\underline{R}}(Qf_*)f^*E\simeq \underline{\underline{R}}(f_*Q)f^*E\simeq f_*\underline{\underline{R}}Qf^*E,
\]
since $f$ is affine. We deduce that
\[
 R^nQ(f_*f^*E)\simeq f_*R^nQ(f^*E)
\quad\text{for every }n\geq0,
\]
whence, by the affine case,
\[
 R^nQ(f_*f^*E)=0\quad\text{for }n>N.
\]
Since $C^\bullet(E)$ has length $N'$, we finally have
\[
 R^nQ(E)=0\quad\text{for }n>N+N',
\]
and we are done.

2°) $S$ is \uline{noetherian}. The components of $C^\bullet(E)$ are sums of terms of the form $f_*f^*(E)$, where $f:U\longrightarrow S$ is the inclusion of a quasi-compact and \uline{separated} open subset. As in case 1°), we have
%%%%%%%%%%  scan idx 206  |  folio 193  %%%%%%%%%%
\SGAPrintPage{193}%
\begin{equation}\tag{*}
 \underline{\underline{R}}(Qf_*)f^*E\simeq \underline{\underline{R}}Q\,f_*f^*E\simeq \underline{\underline{R}}(f_*Q)f^*E.
\end{equation}
Since $U$ is noetherian, $\Qcoh(U)$ has enough objects that are injective in $\Mod(U)$ ([H] II 7.18), so the right derived functor of $f_*:\Qcoh(U)\longrightarrow\Qcoh(S)$ is \textquotedblleft induced\textquotedblright{} by $\underline{\underline{R}}f_*:\mathrm D^+(U)\longrightarrow\mathrm D^+(S)$ when we identify $\mathrm D^+(\Qcoh(U))$ (resp. $\mathrm D^+(\Qcoh(S))$) with a full subcategory of $\mathrm D^+(U)$ (resp. $\mathrm D^+(S)$) by (3.5). Consequently, we have
\[
 \underline{\underline{R}}(f_*Q)f^*E\simeq \underline{\underline{R}}f_*\underline{\underline{R}}Qf^*E.
\]
We deduce a spectral sequence
\[
 E_2^{pq}=R^pf_*R^qQf^*E\Longrightarrow R^{*}(f_*Q)f^*E.
\]
By case 1°), there exists an integer $n(U)$ depending only on $U$ such that $Q$ has cohomological dimension $\leq n(U)$ on $\Mod(U)$. Thus we have
\[
 E_2^{pq}=0\quad\text{for }q>n(U).
\]
But we have
\[
 E_2^{pq}=0\quad\text{for }p>N,
\]
because of the \textquotedblleft cohomological dimension\textquotedblright{} hypothesis on $S$. We therefore have
\[
 R^n(f_*Q)f^*E=0\quad\text{for }n>N+N'',
\]
where $N''$ is an integer greater than all the $n(U)$, with $U$ ranging over the set of intersections $U_{i_0}\cap\ldots\cap U_{i_p}$. Taking account of (*), we obtain
\[
 R^nQ(f_*f^*E)=0\quad\text{for }n>N+N'',
\]
whence finally
\[
 R^nQ(E)=0\quad\text{for }n>N+N'+N'',
\]
%%%%%%%%%%  scan idx 207  |  folio 194  %%%%%%%%%%
\SGAPrintPage{194}%
which completes the proof of a).

\uline{Proof of} b). The coherer, having finite cohomological dimension by a), admits ([H] I 5.3) a right derived functor
\[
 \underline{\underline{R}}Q:\mathrm D(S)\longrightarrow\mathrm D(\Qcoh(S))
\]
(whose restriction to $\mathrm D^+(S)$ coincides with the functor $\underline{\underline{R}}Q$ defined above). For $E\in\Ob\mathrm D(\Qcoh(S))$, there is a canonical functorial isomorphism
\sgaStableTarget{sga6.e02.eq.1.occ2}\begin{equation}\tag{1}
 E\xrightarrow{\sim}\underline{\underline{R}}Q(E),
\end{equation}
since, by (3.5.1), the $E^i$ are acyclic for $Q$ and satisfy $E^i\xrightarrow{\sim}Q(E^i)$. On the other hand, for $E\in\Ob\mathrm D(S)$, we have a canonical functorial homomorphism
\sgaStableTarget{sga6.e02.eq.2.occ2}\begin{equation}\tag{2}
 \underline{\underline{R}}Q(E)\longrightarrow E
\end{equation}
(coming, in the case where the $E^i$ are $Q$-acyclic, from the ordinary adjunction homomorphism between $\phi$ and $Q$). We check trivially that the homomorphisms (1) and (2) make $\underline{\underline{R}}Q:\mathrm D(S)\longrightarrow\mathrm D(\Qcoh(S))$ a right adjoint to $\phi:\mathrm D(\Qcoh(S))\longrightarrow\mathrm D(S)$. The fact that (1) is an isomorphism implies that $\phi$ is fully faithful. Finally, using the spectral sequence
\[
 R^pQ(H^q(E))\Longrightarrow R^{*}Q(E),
\]
(which converges since $Q$ has finite cohomological dimension), we see that (2) is an isomorphism for $E\in\Ob\mathrm D(S)_{\qcoh}$, so $\phi$ has as essential image $\mathrm D(S)_{\qcoh}$. This proves b) and completes the proof of (3.7).

